\chapter{Estado del Arte}
\label{ch:state_of_the_art}

\section{Centros de transformación en redes inteligentes}
En el ecosistema de la Infraestructura de Medición Avanzada (AMI), los centros de transformación de media a baja tensión (MT/BT) representan el nodo neurálgico para la centralización del tráfico de datos. A diferencia de los modelos de distribución dispersos, el modelo masivo europeo y asiático agrupa densidades críticas de abonados (típicamente entre 10 y 100 medidores) bajo el área de cobertura de una única estación transformadora secundaria, justificando técnica y económicamente el despliegue de Unidades Concentradoras de Datos (DCU) fijas \cite{6925952}. El concentrador de datos opera como una pasarela híbrida o \textit{gateway}, gestionando el tráfico descendente hacia los contadores inteligentes a través del propio cableado eléctrico como canal físico mediante tecnologías de comunicación por línea eléctrica de banda estrecha (NB-PLC), y derivando el tráfico ascendente hacia los sistemas centrales de gestión de la distribuidora mediante redes celulares o de fibra óptica \cite{SHARMA2017704, 10843716}. No obstante, la subestación secundaria constituye un entorno hostil para la propagación de alta frecuencia debido a las variaciones dinámicas de la impedancia de acceso, los fenómenos de desadaptación de líneas y la inyección continua de interferencias electromagnéticas \cite{9115397, 9881962}.

\subsection{Esquema de equipos y estrategias de acoplamiento físico}
La inyección de señales NB-PLC en las barras de baja tensión del transformador exige el diseño de Analog Front-Ends (AFE) e interfaces de acoplamiento capacitivo o inductivo capaces de aislar los circuitos digitales de control frente a las tensiones nominales y transitorios de la red (Categoría de Sobretensión IV). Con el fin de superar la atenuación severa del canal eléctrico, las estrategias de inyección física en el centro de transformación se segmentan en arquitecturas monofásicas y trifásicas, teniendo un impacto directo sobre la simetría del canal y el acoplamiento cruzado entre fases (\textit{cross-phase coupling}):
\begin{itemize}
    \item \textbf{Inyección Monofásica:} Presenta ventajas en coste y simplicidad de hardware, pero depende intrínsecamente del acoplamiento natural e indirecto que ocurre en los devanados del transformador y las capacitancias parásitas de las líneas. Bajo escenarios de cargas trifásicas desequilibradas o alta atenuación inductiva, esta estrategia provoca asimetrías severas en la relación señal-ruido (SNR) de las fases no inyectadas.
    \item \textbf{Inyección Trifásica Dirigida:} Requiere un diseño de hardware de alta densidad donde el concentrador gobierna etapas de conmutación activa y filtros sintonizados selectivos para inyectar la señal OFDM de forma síncrona en las tres fases y el neutro \cite{5764420}. Esta aproximación minimiza la presencia de muescas de atenuación (\textit{notches}) en el espectro de la banda CENELEC-A y maximiza la disponibilidad espacial de la subred AMI.
\end{itemize}

El límite operativo crítico en este esquema de hardware no reside en la atenuación resistiva de las líneas conductores, sino en la coexistencia con el ruido conductivo generado por la electrónica de potencia. Mediciones de campo demuestran que las emisiones de alta frecuencia procedentes de fuentes conmutadas (SMPS), inversores fotovoltaicos y convertidores de carga de vehículos eléctricos actúan de forma síncrona en la banda CENELEC-A \cite{10104173}. Este ruido electromagnético satura los amplificadores de bajo ruido (LNA) del receptor del concentrador, degradando la capacidad de corrección de errores de la capa física (PHY) y aislando la cabecera. Esto obliga al desarrollo de procesadores embebidos potentes en el concentrador capaces de soportar esquemas de codificación avanzados y algoritmos de corrección de errores más robustos que los convolucionales estándar, tales como la codificación basada en códigos Turbo o LDPC \cite{9441595,9820886}.

\subsection{Red PRIME y la topología dinámica del Nodo Base}
El protocolo \textit{PoweRline Intelligent Metering Evolution} (PRIME) organiza la subred del centro de transformación mediante una topología en árbol lógica de carácter dinámico (figura \ref{fig:prime_topology}) y \textit{plug-and-play}. Esta arquitectura es gobernada de manera centralizada por el \textbf{Nodo Base} (Base Node), el cual se encuentra embebido en el núcleo de procesamiento del concentrador de datos, mientras que los contadores inteligentes de los usuarios finales actúan inicialmente como \textbf{Nodos de Servicio} (Service Nodes).

\begin{figure}[ht]
    \centering
    % \includegraphics[width=0.8\textwidth]{tu_diagrama_bloques.png}
    \caption{Topología dinámica de red PRIME con niveles de conmutación (\textit{Switches}) coordinados por el Nodo Base.}
    \label{fig:prime_topology}
\end{figure}

Las simulaciones de topología de red NB-PLC demuestran que la estructura ramificada del cableado de baja tensión eleva la tasa de colisiones de tramas en los buffers del \textit{switch} del concentrador \cite{8704772}. Para mitigar este problema, el Nodo Base ejecuta políticas dinámicas en la capa de control de acceso al medio (MAC) y enrutamiento adaptativo OFDM/OFDMA con asignación inteligente de sub-bandas carrier \cite{8362574}:
\begin{enumerate}
    \item \textbf{Promoción a repetidores lógicos (Switches):} Cuando un Nodo de Servicio es incapaz de enlazar directamente con el Nodo Base debido a una baja atenuación o a un foco de ruido local en su feeder, transmite ráfagas de sincronización. El Nodo Base evalúa los nodos candidatos del entorno y promociona a Nodos de Servicio óptimos al rango de \textbf{Switches} (repetidores de señal), regenerando las balizas (\textit{beacons}) para restaurar la conectividad \cite{8693361}.
    \item \textbf{Niveles de Conmutación y Degradación por Latencia:} A pesar de que la inclusión de múltiples niveles de conmutación en cascada amplía la cobertura geográfica, el análisis empírico del rendimiento de la capa de aplicación mediante la pila DLMS/COSEM evidencia que superar los 4 o 5 saltos lógicos incrementa exponencialmente el retardo de propagación \cite{7778783}. Esto penaliza el tiempo medio de lectura masiva de curvas de carga (\textit{load profiles}) y satura el ancho de banda del concentrador de datos.
\end{enumerate}

Asimismo, ante redes complejas y de gran escala donde se manejan masivamente datos multidimensionales de facturación, calidad de suministro y fasores síncronos, las DCU de nueva generación deben integrar en su software módulos avanzados de agregación de datos con preservación de ciberseguridad/privacidad homomórfica y técnicas de poda de datos (\textit{data pruning}) para reducir el impacto del tráfico redundante hacia la cabecera sin alterar la estimación de estado de la red inteligente \cite{9313012, 10268020}. Esta complejidad justifica la transición hacia arquitecturas embebidas multifrecuencia y bandas extendidas capaces de operar de forma flexible más allá de los límites tradicionales de la última milla europea \cite{7007657}.